\section{Introduction}
\label{sec:intro}
Large language model (LLM) agents are increasingly deployed to perform consequential tasks autonomously: booking travel, managing files, executing code, and interacting with APIs on behalf of users~\citep{liu2025agentbench,zhou2024webarena,xu2025theagentcompanybench}. As these agents acquire the ability to commit irreversible actions in real environments, the question becomes: do they know when \emph{not} to act?

Consider a travel-management agent asked to ``cancel the outbound flight and rebook a window seat on the next departure.'' A well-calibrated agent should verify rebooking availability \emph{before} canceling the existing reservation. Yet in practice, we observe that frontier models cancel the flight first, discover that no suitable alternative exists, and then inform the user that the rebooking failed, leaving the user stranded with neither the original nor a new reservation. This failure mode combines irreversible side effects with a belated verbal acknowledgment of the problem. It is distinct from, and more dangerous than, a simple wrong answer in a QA setting.

Abstention, the ability to recognize when a response or action is inappropriate and to refrain accordingly, has been studied extensively for LLMs in question-answering. AbstentionBench~\citep{kirichenko2025abstentionbench} demonstrates that even frontier reasoning models struggle on unanswerable questions. Calibration and uncertainty research~\citep{kapoor2025largelanguagemodelstaught,tomani2024uncertaintybasedabstentionllmsimproves} shows that prompting alone is insufficient and that abstention must often be explicitly taught. On the flip side, over-refusal benchmarks~\citep{cui2025orbench,rottger2024xstest,zhang2025falsereject} establish that the act/abstain boundary is a two-sided trade-off.

However, this line of work, including recent extensions to embodied question answering~\citep{wu2025abstaineqa}, evaluates an answer-or-abstain decision over a single response: the decision cannot alter any environment state, so the worst case is a wrong answer, and there is no tool-call trace against which restraint can be verified. Tool-using agents face a fundamentally different challenge: their abstention decision is \emph{sequential} and \emph{grounded in observable environment state}, requiring multiple tool calls before feasibility can be determined. Safety benchmarks~\citep{andriushchenko2025agentharm,ruan2024identifyingriskslmagents,zhang2025agentsafetybench} measure refusal of explicitly malicious goals but do not address when a well-intentioned task should be abandoned mid-execution. Tool-use reliability benchmarks~\citep{patil2025bfcl,zhang2024toolbehonest,yao2024taubench} test whether agents call tools correctly, not whether they recognize when \emph{no tool call is appropriate}. Clarification-seeking work~\citep{wang2026learningaskllmagents,zhao2026askaskbenchrubricguidedrlvr} focuses on the pre-execution decision but cannot surface failures that emerge mid-trajectory.

To address this gap, we present the first systematic evaluation framework for \emph{agentic abstention}\footnote{Concurrent work \citep{luo2026agentic} also studies agentic abstention by adapting existing benchmarks; see \S\ref{sec:related} for comparison.}: the calibrated ability of tool-using LLM agents to recognize when not to act, spanning both pre-execution reasoning and runtime discovery. The benchmark rests on three design principles:
\textbf{(1) Paired design for calibrated measurement.} Every instance is a pair $(\mathcal{T}^+, \mathcal{T}^-)$ sharing the same sandbox but differing in a single controlled perturbation~$\delta$. The should-act variant~$\mathcal{T}^+$ is fully specified; the should-abstain variant~$\mathcal{T}^-$ introduces exactly one trigger that makes execution incorrect, harmful, or epistemically unjustified. This measures over-abstention and under-abstention all together. By construction, no always-act or always-refuse policy can exceed 50\% Paired Accuracy.
\textbf{(2) Agent-native phase coverage.} We define 8 abstention categories organized by when the trigger becomes observable. Pre-execution categories test recognition from instruction and tool descriptions alone; runtime categories require discovering triggers through tool interaction, testing whether agents can halt rather than push forward after encountering new information, a regime with no analogue in text-only QA abstention.
\textbf{(3) Executable environments with commit-level evaluation.} All 263 paired tasks execute in 42 sandboxed MCP environments with tools categorized as lookup (read-only), verify (validation gates), or commit (state-mutating). Evaluation combines a deterministic commit check on the tool-call trace with an LLM judge on the terminal response, catching failure modes invisible to either signal alone.

To scale evaluation beyond manually authored tasks, we propose \pipeline{}, a fully automated pipeline that synthesizes sandbox environments and generates paired task instances end-to-end. Tasks are validated via deterministic DAG replay at construction time and by cross-family LLM critics. To confirm final task quality, three independent human annotators evaluated a stratified sample of 100 paired tasks, reporting 94--98\% validity per annotator with high inter-annotator agreement.

Evaluating 17 frontier LLMs across 4 agent harnesses, we find that agentic abstention remains unsolved. The best model (Gemini 3.1 Pro) achieves only 59.5\% Paired Accuracy\footnote{The fraction of pairs whose should-act and should-abstain variants are both answered correctly; defined in \S\ref{sec:metrics}.}, and 13 of 17 models stay below 50\%. We also surface failure patterns unique to agentic settings, notably \emph{post-hoc abstention} where agents first commit irreversible actions, and then claim refusal. Abstention capability scales largely independently of task-solving capability, indicating that model scale alone does not resolve the gap.

\textbf{Contributions.} As the first systematic evaluation framework for agentic abstention, this paper makes the following contributions: (1) We present \textbf{the paired-task benchmark \benchmark{}}, built on the first agent-native abstention taxonomy with 8 scenarios organized by when the trigger becomes observable (pre-execution vs.\ runtime) and where it resides (query, environment state, or tools). Its should-act and should-abstain paired design measures over- and under-abstention all together, with 263 paired tasks instantiated in 42 executable MCP sandbox environments and a composite metric (commit check and LLM judge) that isolates failure modes invisible to either signal alone. (2) We propose \textbf{a fully automated pipeline \pipeline{}} that synthesizes sandbox environments and generates paired tasks end-to-end, validated by deterministic DAG replay and cross-family LLM validation; on-demand regeneration mitigates data contamination as models evolve. (3) We identify \textbf{failure modes unique to agentic abstention}: the best of 17 frontier models reaches only 59.5\% Paired Accuracy, and abstention capability scales independently of task-solving capability.
